% Part: first-order-logic
% Chapter: completeness
% Section: completeness-theorem

\documentclass[../../../include/open-logic-section]{subfiles}

\begin{document}

\iftag{FOL}
      {\olfileid[id]{fol}{com}{cth}}
      {\olfileid[id]{pl}{com}{cth}}

\olsection{Teorema Kelengkapan}

\begin{explain}
Mari kita gabungkan hasil-hasil yang telah diperoleh: kita sampai pada
teorema kelengkapan.
\end{explain}

\begin{thm}[Teorema Kelengkapan]
\ollabel{thm:completeness}
Misalkan $\Gamma$ adalah suatu himpunan !!{sentence}. Jika $\Gamma$
konsisten, maka himpunan itu dapat dipenuhi.
\end{thm}

\begin{proof}
Andaikan $\Gamma$ konsisten.\iftag{FOL}{ Berdasarkan
  \olref[hen]{lem:henkin}, terdapat himpunan konsisten dan jenuh
  $\Gamma' \supseteq \Gamma$.}{} Berdasarkan
\olref[lin]{lem:lindenbaum}, terdapat $\Gamma^* \supseteq
\iftag{FOL}{\Gamma'}{\Gamma}$ yang konsisten dan !!{complete}.
\iftag{FOL}{
  Karena $\Gamma' \subseteq \Gamma^*$, untuk setiap
  !!{formula}~$!A(x)$, $\Gamma^*$ memuat !!a{sentence} berbentuk
  \iftag{prvEx}
        {$\lexists[x][!A(x)] \lif !A(c)$}
        {$\lnot\lforall[x][!A(x)] \lif \lnot !A(c)$}
  sehingga $\Gamma^*$ jenuh. Jika $\Gamma$ tidak memuat~$\eq$, maka
  berdasarkan}{Berdasarkan}
\olref[mod]{lem:truth},
$\iftag{FOL}{\Sat{M(\Gamma^*)}}{\pSat{v(\Gamma^*)}}{!A}$ jika dan hanya
jika $!A \in \Gamma^*$. Khususnya, dari sini diperoleh bahwa untuk
setiap $!A \in \Gamma$,
$\iftag{FOL}{\Sat{M(\Gamma^*)}}{\pSat{v(\Gamma^*)}}{!A}$; jadi $\Gamma$
dapat dipenuhi.\iftag{FOL}{
  Jika $\Gamma$ memuat~$\eq$, maka berdasarkan
  \olref[ide]{lem:truth}, untuk setiap !!{sentence}~$!A$,
  $\Sat{\equivclass{M}{\approx}}{!A}$ jika dan hanya jika $!A \in
  \Gamma^*$. Khususnya, $\Sat{\equivclass{M}{\approx}}{!A}$ untuk setiap
  $!A \in \Gamma$; jadi $\Gamma$ dapat dipenuhi.}{}
\end{proof}

\begin{cor}[Teorema Kelengkapan, Versi Kedua]
\ollabel{cor:completeness}
Untuk setiap $\Gamma$ dan !!{sentence}~$!A$: jika $\Gamma \Entails !A$,
maka $\Gamma \Proves !A$.
\end{cor}

\begin{proof}
Perhatikan bahwa $\Gamma$ dalam \olref{cor:completeness} dan
\olref{thm:completeness} dikuantifikasi secara universal. Untuk
menghindari kekeliruan, mari kita nyatakan kembali
\olref{thm:completeness} dengan variabel lain: untuk sebarang himpunan
!!{sentence}~$\Delta$, jika $\Delta$ konsisten, maka himpunan itu dapat
dipenuhi. Dengan mengambil kontraposisinya, jika $\Delta$ tidak dapat
dipenuhi, maka $\Delta$ tak konsisten. Kita akan menggunakan fakta ini
untuk membuktikan korolari tersebut.

Andaikan $\Gamma \Entails !A$. Maka $\Gamma \cup \{\lnot !A\}$ tidak
dapat dipenuhi berdasarkan \olref[syn][sem]{prop:entails-unsat}. Dengan
mengambil $\Gamma \cup \{\lnot !A\}$ sebagai $\Delta$, versi sebelumnya
dari \olref{thm:completeness} menyatakan bahwa $\Gamma \cup \{\lnot
!A\}$ tak konsisten. Berdasarkan
\iftag{FOL}{%
  \tagrefs{prfAX/{fol:axd:prv:prop:prov-incons},
    prfSC/{fol:seq:prv:prop:prov-incons},
    prfND/{fol:ntd:prv:prop:prov-incons},
    prfTab/{fol:tab:prv:prop:prov-incons}}}{%
  \tagrefs{prfAX/{pl:axd:prv:prop:prov-incons},
    prfSC/{pl:seq:prv:prop:prov-incons},
    prfND/{pl:ntd:prv:prop:prov-incons},
    prfTab/{pl:tab:prv:prop:prov-incons}}},
$\Gamma \Proves !A$.
\end{proof}

\tagprob{FOL}
\begin{prob}
Gunakan \olref[fol][com][cth]{cor:completeness} untuk membuktikan
\olref[fol][com][cth]{thm:completeness}, sehingga ditunjukkan bahwa kedua
rumusan teorema kelengkapan tersebut ekuivalen.
\end{prob}
\tagendprob

\tagprob{notFOL}
\begin{prob}
Gunakan \olref[pl][com][cth]{cor:completeness} untuk membuktikan
\olref[pl][com][cth]{thm:completeness}, sehingga ditunjukkan bahwa kedua
rumusan teorema kelengkapan tersebut ekuivalen.
\end{prob}
\tagendprob

\tagprob{FOL}
\begin{prob}
Agar suatu sistem !!a{derivation} lengkap, aturan-aturannya harus cukup
kuat untuk membuktikan bahwa setiap himpunan yang tidak dapat dipenuhi
tak konsisten. Aturan !!{derivation} manakah yang diperlukan untuk
membuktikan kelengkapan? Adakah aturan-aturan ini yang tidak digunakan
di mana pun dalam pembuktian? Untuk menjawab pertanyaan-pertanyaan
tersebut, buatlah daftar atau diagram yang menunjukkan aturan
!!{derivation} mana yang digunakan dalam hasil mana saja yang mengarah pada
pembuktian \olref[fol][com][cth]{thm:completeness}. Pastikan Anda mencatat
setiap penggunaan aturan yang tersirat dalam pembuktian-pembuktian ini.
\end{prob}
\tagendprob

\tagprob{notFOL}
\begin{prob}
Agar suatu sistem !!a{derivation} lengkap, aturan-aturannya harus cukup
kuat untuk membuktikan bahwa setiap himpunan yang tidak dapat dipenuhi
tak konsisten. Aturan !!{derivation} manakah yang diperlukan untuk
membuktikan kelengkapan? Adakah aturan-aturan ini yang tidak digunakan
di mana pun dalam pembuktian? Untuk menjawab pertanyaan-pertanyaan
tersebut, buatlah daftar atau diagram yang menunjukkan aturan
!!{derivation} mana yang digunakan dalam hasil mana saja yang mengarah pada
pembuktian \olref[pl][com][cth]{thm:completeness}. Pastikan Anda mencatat
setiap penggunaan aturan yang tersirat dalam pembuktian-pembuktian ini.
\end{prob}
\tagendprob

\end{document}
